\section{Discussion and Limitations}
\label{sec:discussion}
% Budget: 0.5 page. Source: DECISIONS.md s16, s25.5, s25.6, s26.3, s26.5, s14 erratum.

\textbf{When this matters.} The attack needs benign traffic the adversary can reproduce with
near-duplicate fidelity, and CICIDS2017 makes that cheap because its own benign class is narrow
and repetitive (Section~\ref{sec:degeneracy}) — replaying the pool's own statistics with no
perturbation already lands one row in ten within twin-fraction range
(Section~\ref{sec:attack}). Production networks with genuinely diverse, high-entropy benign
traffic would raise the fidelity bar the attacker has to clear; networks that do not — IoT
fleets, ICS traffic, automated service-to-service calls, anything whose "normal" is a small set
of repeated request shapes — are closer to CICIDS2017's regime than to a diverse enterprise
network, and are where this channel should be expected to be cheapest to exploit. We do not
measure this on production traffic; it is a prediction from the mechanism, not a finding.

\textbf{Practical advice.} Cap the honeypot's effective share of the training set; it is the one
defense here that costs nothing to deploy, needs no label information or trusted reference data
beyond the cap value, and is undominated by every alternative we test
(Section~\ref{sec:defenses}). Do not trust a label-consistency check that learns from the loop's
own data to catch this poison — it cannot, because the poison is consistent with the loop's own
labelling policy by construction; a check anchored to data outside the loop can.

\textbf{Limitations.} This is an offline simulation of the loop, not a deployed system: the
adversary's batches are drawn from a fixed pool rather than generated against a live honeypot,
and there is no network-level cost to reaching the decoy in the first place. We measure flow
features only, not payload content, so any attack or defense that depends on packet contents is
out of scope. CICIDS2017 answers the degeneracy measurement and the fidelity-versus-damage curve
but carries no learning claim (Section~\ref{sec:degeneracy}); the synthetic dataset answers the
honest-loop learning claim but has no attacker-cost profile of its own. Five seeds is enough to
establish some effects and not others: XGBoost's fixed-threshold FPR trajectory splits into two
distinct per-seed groups rather than clustering around a mean (Section~\ref{sec:attack}), and
RF's damage at large mimicry distance needed a dedicated 15-seed follow-up to confirm — the
original 5-seed, 36-cell grid family lacked the power to resolve it, even though the effect was
real (Section~\ref{sec:attack}). We only had the budget to run that follow-up for RF; the same
question is open for XGBoost. Further limitations — the partition's file-order versus verified
timestamp, and ShareCap's cap range and adaptive-attacker coverage — are in
Appendix~\ref{sec:appendix}.
